\subsubsection{Instrumento, patrones y condiciones ambientales}

\begin{tabular}{ll}
Marca:      & Mitutoyo\\[3pt]
Código:     & 530-115\\[3pt]
Serie:      & 2013G454\\[3pt]
Resolución: & $0.05\text{ mm}$\\[3pt]
\end{tabular}

\vspace{6pt}
Procedimiento por comparación con instrumento de referencia según norma\\ PC-012 INACAL (\textit{Procedimiento de calibración de pie de rey}).

\vspace{12pt}
\subsubsection*{Equipos e instrumentos utilizados:}
\begin{tabular}{ll}
Bloques patrón de longitud & \\[3pt]
Cilindro patrón            & \\[3pt]
Anillo patrón              & \\[3pt]
Mesa de planitud           & \\[3pt]
Termómetro de contacto     & Exactitud $\pm 0.5\ ^\circ\text{C}$
\end{tabular}

\vspace{12pt}
\subsubsection*{Nomenclatura:}
\begin{tabular}{ll}
BU  & Bloques Utilizados\\[3pt]
VP  & Valor Patrón\\[3pt]
VPC & Valor Patrón Corregido\\[3pt]
PC  & Promedio Corregido\\[3pt]
$X_j$ & Lectura del pie de rey en la posición $j$ (superior, central, inferior)\\[3pt]
$D_j$ & Diferencia entre la lectura $j$ y el valor patrón corregido\\[3pt]
$E_k$ & Dispersión de las tres lecturas en el punto de calibración $k$\\[3pt]
$S_{E-I}$ & Error por cambio de la escala de exteriores a interiores\\[3pt]
$S_{E-P}$ & Error por cambio de la escala de exteriores a profundidad
\end{tabular}

\subsubsection{Datos experimentales}
\label{sec:pie-rey-datos}

Las mediciones se realizaron en cinco puntos distribuidos sobre el alcance del
instrumento, formando la pila de bloques patrón por acumulación: $50$ mm (un
bloque), $100$ mm (un bloque), $150$ mm (dos bloques), $200$ mm (tres bloques) y
$250$ mm (cuatro bloques). En cada punto se efectuaron tres mediciones
correspondientes a las posiciones superior, central e inferior de las mordazas de
medición de exteriores, conforme al ítem 10.1.2 de la norma PC-012. La temperatura
registrada fue de $21.5\ ^\circ$C en los cinco puntos, tal como se realizaron las
mediciones.

El pie de rey 530-115 es un instrumento de vernier cuya escala principal está
dividida en milímetros y cuyo nonio consta de 20 divisiones que abarcan 19 mm. La
diferencia entre una división de la escala principal y una del nonio define la
resolución del instrumento, $d = 1 - 19/20 = 0.05$ mm, de modo que toda lectura es
un múltiplo de dicha división. Las lecturas que se presentan a continuación respetan
esa graduación.

La Tabla~\ref{tab:pie-rey-datos} recoge las mediciones de los cinco puntos. La
columna BU indica los bloques que componen cada pila y la columna $VP$ el valor
nominal de la pila resultante.

\begin{table}[H]
\centering
\caption{Mediciones de calibración del pie de rey. $X_1$, $X_2$ y $X_3$ son las lecturas en las posiciones superior, central e inferior de medición de exteriores.}
\label{tab:pie-rey-datos}
\begin{tabular}{lrrrrr}
\toprule
Bloques patrón (mm) & $VP$ (mm) & $X_1$ (mm) & $X_2$ (mm) & $X_3$ (mm) & $T$ ($^\circ$C) \\
\midrule
50 & 50.00 & 50.00 & 50.00 & 50.00 & 21.5 \\
100 & 100.00 & 100.00 & 100.05 & 100.00 & 21.5 \\
50, 100 & 150.00 & 150.00 & 150.05 & 150.05 & 21.5 \\
25, 75, 100 & 200.00 & 200.00 & 200.05 & 200.05 & 21.5 \\
25, 50, 75, 100 & 250.00 & 250.00 & 250.00 & 250.00 & 21.5 \\
\bottomrule
\end{tabular}
\end{table}

